\documentclass[a4paper,12pt]{article}

\usepackage[french]{babel}
\usepackage[utf8]{inputenc}
\usepackage[T1]{fontenc}
\usepackage{graphicx}

\title{Sujet 2}
\author{HAMED Rémy, SOARES Enzo, TURBIEZ Denis}
\date{20 Mars 2022}

\begin{document}

\maketitle
\newpage

\tableofcontents
\newpage

\part{Proposition de sujet}

\paragraph{Maintenance, entretien et vigilance temporaire à domicile}
En cas d'absence temporaire d'une personne de son domicile (pour vacance, hospitalisation, déplacement professionnel …), proposé des personnes pouvant accomplir des tâches de maintenance, d'entretien et de surveillance (légère, pas de vigile, de ronde ou de vidéo surveillance fournit) du domicile afin que la personne soit sûre du bon état de son domicile à son retour.

\newpage
\part{Descriptif fonctionnel}

\section{Client}
Avant ou après son départ de son domicile, le client peut utiliser l'application afin de trouver une personne pouvant assurer les tâches de maintenance, d'entretien et de surveillance du domicile. Pour cela, 2 parcours sont possible.

\subsection{Parcours 1}
Le client cherche dans une liste de prestataire vérifié, selon plusieurs critère de filtre possible, afin de trouver la personne la plus qualifiée pour sa demande.\\
Le client peut ensuite rentrer en contact avec le prestataire (téléphone, mail, messagerie instantané, rendez vous physique) afin de discuter des modalité de la prestation (tâche à effectuer, durée, rémunération …).\\
Le client peut (dois ?) également fournir une liste des tâches à effectuer ainsi que leur fréquence.\\
Une rencontre entre le client et le prestataire peut être organisé avant le début de la période d'utilisation afin de fixer les dernière condition et expliquer le déroulement de certaine actions.\\

\subsection{Parcours 2}
Le client prépare une fiche de domicile comprenant les différentes modalité de la prestation et la poste sur l'application.\\
Les prestataires peuvent postuler pour une offre. Chaque prestataire postulant est mis en contact avec le client, afin d'avoir une évaluation par le client du prestataire.\\
Une fois le choix fait une rencontre entre le client et le prestataire peut être organisé pour fixer les dernières condition et expliquer certaine actions.

\subsection{Après prestation}
Quand la prestation est terminée, le client et le prestataire doivent faire un état des lieux afin de vérifier la bonne tenue de tout les engagements.\\
Une fois ceci fait le paiement du prestataire doit être fait ou prévu (émission d'une facture par l'application).\\
En cas de litige ?

\section{Prestataire}
Le prestataire à un compte sur la plateforme ayant vérifier son statut et son identité.\\
Il peut fournir certaine information relative aux services qu'il offre.\\
Pour entrer en contact avec le client 2 parcours sont possible.

\subsection{Parcours 1}
Le prestataire est choisi par un client pour une prestation n'étant pas encore enregistrée sur la plateforme. Le client et le prestataire sont mis en contact (téléphone, mail, messagerie instantané, rencontre physique …), afin de discuter des modalités de la prestation (tâches, durée, rémunération).\\
Une fois ceci fait, une rencontre physique est organisé afin de fixer les modalité et le déroulement d'actions spécifique.\\

\subsection{Parcours 2}
Le prestataire répond à une offre déjà mise sur l'application présentant donc déjà les modalité. À la candidature le prestataire et le client sont mis en relation afin que le client puisse en apprendre plus sur le prestataire et faire son choix.\\
Quand le prestataire est sélectionné par le client, un rendez vous physique est organisé afin que les 2 partie puissent convenir des dernière modalité de la prestation.

\subsection{Après la prestation}
Une fois la prestation terminée, le client et le prestataire doivent se rencontrer au domicile en question afin de faire un état des lieux.\\
Ceci fait le prestataire doit émettre une facture et si possible se faire payer.\\
En cas de litige ?

\newpage
\part{Traitement client}

\section{Mobile}

\subsection{Client}
Sur un appareil mobile, le client peut voir l'avancement des tâches réalisé dans son domicile, les prestataires ayant répondu à l'offre postée, les conversations qu'il a eu avec les différents prestataire ayant été mis en contact avec lui.

\subsection{Prestataire}
Sur un appareil mobile, le prestataire peut avoir accès à la liste des taches qu'il doit effectuer pour une prestation. Il peut accéder aux offres auxquelles il a répondu ainsi que les conversations qu'il à eu avec les différents client ayant pris contact avec lui pour une prestation.\\

\section{Ordinateur}

\subsection{Client}
Sur un appareil de bureau, le client peut interagir avec toute les fonctionnalité de l'application, comprenant donc\footnote{Liste non exhaustive} : 
\begin{itemize}
\item Recherche de prestataire
\item Sélection de prestataire
\item Publication d'une offre de prestation
\item Mise en place d'une liste de tâche
\item Suivi des tâches effectuer sur une prestation
\item Suivi des conversations avec le prestataire
\end{itemize}

\subsection{Prestataire}
Sur un appareil de bureau, le prestataire peut interagir avec toute les fonctionnalité de l'application, comprenant donc \footnote{Liste non exhaustive}
\begin{itemize}
\item Recherche d'offre
\item Postuler à des offres
\item Proposer ses services (modification des services proposé aux personne)
\item Converser avec les clients des offres auquel il a répondu
\item Converser avec les clients ayant sélectionner leur profil pour une prestation
\item Suivi des tâches à réaliser
\end{itemize}

\newpage
\part{Référent Visuel}

\section{Logo}
\includegraphics{sujet2.png}

\section{Nom de l'entreprise}
Le nom de l'entreprise doit évoqué au client un sentiment de sécurité.\\
C'est pourquoi nous avons choisi comme nom pour notre service secuHouse

\section{Slogan}
Une maison préservée

\end{document}